\uuid{vMno}
\exo7id{4512}
\titre{exo7 4512}
\auteur{quercia}
\organisation{exo7}
\datecreate{2010-03-14}
\isIndication{false}
\isCorrection{true}
\chapitre{Suite et série de fonctions}
\sousChapitre{Convergence simple, uniforme, normale}
\cpge{mpsi-mp2i.nombres-reels-suites.d, pcsi.nombres-reels-suites.c, ptsi.nombres-reels-suites.c, mpsi-mp2i.limites-continuite.b, pcsi.limites-continuite.b, ptsi.limites-continuite.b, mpsi-mp2i.nombres-reels-suites.i, pcsi.nombres-reels-suites.g, ptsi.nombres-reels-suites.g, bcpst1.suites-reelles, bcpst1.limites-continuite.c, mp-mpi.suites-series-fonctions.a, pc.suites-series-fonctions.a, psi.suites-series-fonctions.a, mp-mpi.topologie-espaces-normes.b}
\module{Analyse}
\niveau{L2}
\difficulte{}

\contenu{
\question{
Soit $f : {[-1,1]} \to {[-1,1]}$ une fonction continue vérifiant :
$\forall\ x \ne 0,\quad |f(x)| < |x|$.

On pose $f_0(x) = x$, puis $f_{n+1}(x) = f(f_n(x))$.
\'Etudier la convergence simple des $f_n$.
}
\reponse{
$\bigl(|f_n(x)|\bigr)$ décro\^\i t donc tend vers $L$. On extrait une
sous suite $(f_{\varphi(n)})$ convergeant vers $\ell  \Rightarrow  |\ell| = L$.

La sous-suite $(f_{\varphi(n)+1})$ converge vers
$f(\ell)  \Rightarrow  |f(\ell)| = L  \Rightarrow  L = 0$.
}
}
